%---------------------------------------------------------------------
% Urdu edition, units 5-9.
%---------------------------------------------------------------------

\section*{یونٹ \textenglish{5}: میٹرکس}
\markboth{یونٹ 5: میٹرکس}{}

\begin{nukta}[ابعاد کا میل، اور وہ اصول جو حساب جیسا نہیں]
\(A_{m\times n}B_{n\times p}\) کی ضرب تبھی ممکن ہے جب \emph{اندرونی} ابعاد
برابر ہوں، اور حاصل ضرب \(m\times p\) ہوتا ہے۔

\medskip
جمع اور عددی ضرب عام حساب جیسی ہیں۔ میٹرکس کی ضرب نہیں: عموماً
\(AB\neq BA\)۔ منقلب کے اصول اسی سے نکلتے ہیں:
\[
(A+B)^{t}=A^{t}+B^{t},\qquad (A^{t})^{t}=A,\qquad
\boxed{(AB)^{t}=B^{t}A^{t}}
\]
آخری اصول میں ترتیب \emph{الٹ} جاتی ہے۔ یہی بات معکوس پر بھی لاگو ہے:
\((AB)^{-1}=B^{-1}A^{-1}\)۔
\end{nukta}

\begin{sawal}{سوال \textenglish{14}}
اگر \(A=\mat{1&2&3\\-1&-2&3\\3&4&5}\) اور \(B=\mat{4&-1&1\\1&4&3\\3&-3&-2}\)
ہوں تو \(AB\) نکالیے، اور تصدیق کیجیے کہ \((A+B)^{t}=A^{t}+B^{t}\)۔
\end{sawal}

\begin{hal}
\textbf{\(AB\) کی ضرب۔} \(A\) کی سطر کو \(B\) کے ستون سے ضرب دیجیے۔ پہلی سطر
\((1,\,2,\,3)\) کو \(B\) کے تینوں ستونوں کے ساتھ:
\[
\begin{aligned}
(1)(4)+(2)(1)+(3)(3)&=4+2+9=15\\
(1)(-1)+(2)(4)+(3)(-3)&=-1+8-9=-2\\
(1)(1)+(2)(3)+(3)(-2)&=1+6-6=1
\end{aligned}
\]
اسی طرح دوسری اور تیسری سطر کے لیے:
\[
AB=\mat{15&-2&1\\3&-16&-13\\31&-2&5}
\]

\medskip
\textbf{تصدیق۔}
\[
A+B=\mat{5&1&4\\0&2&6\\6&1&3}
\;\Longrightarrow\;
(A+B)^{t}=\mat{5&0&6\\1&2&1\\4&6&3}
\]
\[
A^{t}+B^{t}=\mat{1&-1&3\\2&-2&4\\3&3&5}+\mat{4&1&3\\-1&4&-3\\1&3&-2}
=\mat{5&0&6\\1&2&1\\4&6&3}
\]
دونوں برابر ہیں۔

\jawab{\(AB=\mat{15&-2&1\\3&-16&-13\\31&-2&5}\)، اور
\((A+B)^{t}=A^{t}+B^{t}\) تصدیق شدہ۔}
\end{hal}

\begin{ghalti}
\((AB)^{t}=B^{t}A^{t}\) ہے، \(A^{t}B^{t}\) نہیں۔ یہ الٹ پھیر کوئی رواج نہیں،
ابعاد کے میل کی مجبوری ہے --- غلط ترتیب سے اکثر تو میٹرکس کی شکل ہی غلط بنتی ہے۔
\end{ghalti}

\begin{sawal}{سوال \textenglish{15}}
برگر بارن کی تین شاخیں فرائز، برگر اور مشروب بیچتی ہیں۔ روزانہ فروخت یہ ہے:
شاخ \textenglish{I} --- \textenglish{500} فرائز، \textenglish{150} برگر،
\textenglish{300} مشروب؛ شاخ \textenglish{II} --- \textenglish{900}،
\textenglish{440}، \textenglish{830}؛ شاخ \textenglish{III} ---
\textenglish{700}، \textenglish{580}، \textenglish{1200}۔ قیمتیں بالترتیب
\textenglish{\$0.92}، \textenglish{\$1.54} اور \textenglish{\$0.58} ہیں۔ تینوں
شاخوں کی کل روزانہ آمدنی نکالیے۔
\end{sawal}

\begin{hal}
فروخت کو \(3\times3\) میٹرکس (سطر = شاخ) اور قیمتوں کو \(3\times1\) ستون میں
لکھیے۔ اندرونی ابعاد \(3=3\) ہیں، اس لیے حاصل ضرب \(3\times1\) --- یعنی ہر شاخ
کی ایک آمدنی:
\[
SP^{t}=\mat{500&150&300\\900&440&830\\700&580&1200}\mat{0.92\\1.54\\0.58}
\]
\[
\begin{aligned}
\text{شاخ I}  &: 460+231+174=865.00\\
\text{شاخ II} &: 828+677.60+481.40=1987.00\\
\text{شاخ III}&: 644+893.20+696=2233.20
\end{aligned}
\]
\[
865.00+1987.00+2233.20=5085.20
\]

\jawab{شاخوں کی آمدنی \textenglish{\$865.00}، \textenglish{\$1{,}987.00} اور
\textenglish{\$2{,}233.20}؛ کل روزانہ آمدنی
\textbf{\textenglish{\$5{,}085.20}}۔}
\end{hal}

\begin{ghalti}
اگر قیمتوں کو \(1\times3\) سطر کے طور پر دیا جائے تو \(SP\) کی ضرب ممکن ہی نہیں۔
کاپی پر لکھیے: ''\(SP\) کے ابعاد میل نہیں کھاتے، اس لیے \(SP^{t}\) نکالا گیا''۔
ممتحن اسی مشاہدے کے نمبر دیتا ہے۔
\end{ghalti}

%---------------------------------------------------------------------
\section*{یونٹ \textenglish{6}: محدد اور معکوس}
\markboth{یونٹ 6: محدد اور معکوس}{}

\begin{nukta}[ہم عامل، معکوس، کریمر]
ہم عامل \(C_{ij}=(-1)^{i+j}M_{ij}\)، جہاں \(M_{ij}\) وہ محدد ہے جو سطر \(i\)
اور ستون \(j\) حذف کرنے کے بعد بچے۔ علامتوں کا نقشہ:
\[
\mat{+&-&+\\-&+&-\\+&-&+}
\]
\[
\det A=\sum_j a_{ij}C_{ij},\qquad
\operatorname{adj}A=\bigl[C_{ij}\bigr]^{t},\qquad
A^{-1}=\frac{1}{\det A}\operatorname{adj}A
\]
\textbf{کریمر کا اصول:} \(Ax=b\) کے لیے \(x_k=\dfrac{D_k}{D}\)، جہاں
\(D=\det A\) اور \(D_k\) میں \(A\) کا \(k\)واں ستون \(b\) سے بدل دیا جاتا ہے۔
شرط یہ ہے کہ \(D\neq 0\) ہو۔
\end{nukta}

\begin{nukta}[محدد کے چار مختصر راستے]
\begin{itemize}
  \item اس سطر یا ستون سے پھیلائیے جس میں سب سے زیادہ صفر ہوں۔
  \item پوری سطر یا ستون کا مشترک عامل باہر نکل آتا ہے۔
  \item دو سطریں (یا ستون) \emph{یکساں} یا متناسب ہوں تو \(\det=0\)۔
  \item ایک سطر کا مضاعف دوسری میں جمع کرنے سے محدد نہیں بدلتا۔
\end{itemize}
\end{nukta}

\begin{sawal}{سوال \textenglish{16}}
محدد نکالیے: \quad
\textbf{(الف)} \(\dt{2&3&-1\\1&1&0\\2&-3&5}\) \qquad
\textbf{(ب)} \(\dt{2a&a&a\\b&2b&b\\c&c&2c}\)
\end{sawal}

\begin{hal}
\textbf{(الف)} دوسری سطر میں ایک صفر ہے، اسی سے پھیلائیے:
\[
2\dt{1&0\\-3&5}-3\dt{1&0\\2&5}+(-1)\dt{1&1\\2&-3}
=2(5)-3(5)-1(-5)=10-15+5=0
\]
محدد کا صفر ہونا بتاتا ہے کہ سطریں آپس میں وابستہ ہیں۔

\medskip
\textbf{(ب)} پہلی سطر سے \(a\)، دوسری سے \(b\) اور تیسری سے \(c\) باہر نکالیے:
\[
\dt{2a&a&a\\b&2b&b\\c&c&2c}=abc\dt{2&1&1\\1&2&1\\1&1&2}
\]
\[
\dt{2&1&1\\1&2&1\\1&1&2}=2(4-1)-1(2-1)+1(1-2)=6-1-1=4
\]
\[
\text{قدر}=4abc
\]

\jawab{(الف) \(0\) \quad (ب) \(4abc\)}
\end{hal}

\begin{sawal}{سوال \textenglish{17}}
\(A=\mat{1&2&3\\1&3&5\\1&5&12}\) کا معکوس ہم عاملوں کے طریقے سے نکالیے۔
\end{sawal}

\begin{hal}
پہلے محدد:
\[
\det A=1(36-25)-2(12-5)+3(5-3)=11-14+6=3 \neq 0
\]
معکوس موجود ہے۔ ہم عامل نکال کر \emph{منقلب} لیجیے (یہی ملحقہ ہے):
\[
\operatorname{adj}A=\mat{11&-9&1\\-7&9&-2\\2&-3&1}
\]
\[
A^{-1}=\frac{1}{3}\mat{11&-9&1\\-7&9&-2\\2&-3&1}
\]
\parakh{\(A\) کی پہلی سطر \((1,2,3)\) کو \(A^{-1}\) کے پہلے ستون
\(\tfrac13(11,-7,2)^{t}\) سے ضرب دیجیے:
\(\tfrac{11-14+6}{3}=1\)۔ تیسری سطر \((1,5,12)\) کو اسی ستون سے:
\(\tfrac{11-35+24}{3}=0\)۔ درست ہے۔}

\jawab{\(A^{-1}=\dfrac{1}{3}\mat{11&-9&1\\-7&9&-2\\2&-3&1}\)}
\end{hal}

\begin{ghalti}
ملحقہ ہم عاملوں کے میٹرکس کا \textbf{منقلب} ہے۔ منقلب بھول جانے سے
\(AA^{-1}=I\) خاص طور پر \emph{غیر قطری} خانوں میں ٹوٹتا ہے --- اسی لیے اوپر
پڑتال میں ایک غیر قطری خانہ بھی جانچا گیا ہے، صرف قطری نہیں۔
\end{ghalti}

\begin{sawal}{سوال \textenglish{18}}
کریمر کے اصول سے حل کیجیے:
\[
2x+y+3z=3,\qquad x+y-2z=0,\qquad -3x-y+2z=-4
\]
\end{sawal}

\begin{hal}
\[
D=\dt{2&1&3\\1&1&-2\\-3&-1&2}=2(2-2)-1(2-6)+3(-1+3)=0+4+6=10
\]
\(D\neq 0\)، اس لیے کریمر کا اصول لاگو ہوتا ہے۔ ایک ایک ستون کو دائیں طرف کے
\((3,0,-4)^{t}\) سے بدلیے:
\[
D_x=\dt{3&1&3\\0&1&-2\\-4&-1&2}=20,\quad
D_y=\dt{2&3&3\\1&0&-2\\-3&-4&2}=-16,\quad
D_z=\dt{2&1&3\\1&1&0\\-3&-1&-4}=2
\]
\[
x=\frac{20}{10}=2,\qquad
y=\frac{-16}{10}=-\frac85,\qquad
z=\frac{2}{10}=\frac15
\]
\parakh{پہلی مساوات: \(4-\tfrac85+\tfrac35=4-1=3\)۔
دوسری: \(2-\tfrac85-\tfrac25=2-2=0\)۔
تیسری: \(-6+\tfrac85+\tfrac25=-6+2=-4\)۔ تینوں پوری ہوتی ہیں۔}

\jawab{\(x=2\)، \(y=-\tfrac85\)، \(z=\tfrac15\)}
\end{hal}

\begin{sawal}{سوال \textenglish{19}}
کریمر کے اصول سے حل کیجیے:
\[
x+y=2,\qquad 2x-z=1,\qquad 2y-3z=-1
\]
\end{sawal}

\begin{hal}
سب سے پہلے غائب متغیرات کے صفر عامل \emph{لکھ} لیجیے، ورنہ اعداد غلط ستون میں
چلے جائیں گے:
\[
x+y+0z=2,\qquad 2x+0y-z=1,\qquad 0x+2y-3z=-1
\]
\[
D=\dt{1&1&0\\2&0&-1\\0&2&-3}=1(0+2)-1(-6-0)+0=2+6=8
\]
\[
D_x=\dt{2&1&0\\1&0&-1\\-1&2&-3}=8,\quad
D_y=\dt{1&2&0\\2&1&-1\\0&-1&-3}=8,\quad
D_z=\dt{1&1&2\\2&0&1\\0&2&-1}=8
\]
\[
x=\frac88=1,\qquad y=\frac88=1,\qquad z=\frac88=1
\]
\parakh{\(1+1=2\)؛ \(2(1)-1=1\)؛ \(2(1)-3(1)=-1\)۔ تینوں درست۔}

\jawab{\(x=1\)، \(y=1\)، \(z=1\)}
\end{hal}

\begin{ghalti}
کریمر لگانے سے \emph{پہلے} ہمیشہ \(D\) نکالیے۔ اگر \(D=0\) ہو تو یہ اصول لاگو
ہی نہیں ہوتا، اور درست جواب اعداد نہیں بلکہ نظام کے بارے میں ایک بیان ہوتا ہے:
''کوئی حل نہیں'' یا ''لامتناہی حل''۔ صفر محدد والے نظام سے اعداد نکل آنا اس بات
کی علامت ہے کہ کہیں حساب کی غلطی ہوئی ہے۔
\end{ghalti}

%---------------------------------------------------------------------
\section*{یونٹ \textenglish{7}: حد اور تفرق}
\markboth{یونٹ 7: حد اور تفرق}{}

\begin{nukta}[حد کی تین صورتیں، اور ہر ایک کا علاج]
\begin{itemize}
  \item \textbf{تعویض چل جائے:} اگر \(f\) نقطہ \(a\) پر متصل ہو تو
        \(\lim_{x\to a}f(x)=f(a)\)۔ ہر بار پہلے یہی آزمائیے۔
  \item \textbf{\(\tfrac00\) کی صورت:} صفر بننے والے عامل منسوخ ہوتے ہیں۔
        تحلیل کیجیے، پھر دوبارہ تعویض۔
  \item \textbf{\(x\to\infty\):} شمار کنندہ اور نسب نما دونوں کو نسب نما کی سب
        سے بڑی قوت سے تقسیم کیجیے، پھر \(\tfrac{1}{x^{n}}\to 0\) لگائیے۔
\end{itemize}
تفرق کے اصول:
\[
\dd{}{x}x^{n}=nx^{n-1},\quad
(uv)'=u'v+uv',\quad
\left(\frac uv\right)'=\frac{u'v-uv'}{v^{2}},\quad
\bigl(f(g(x))\bigr)'=f'(g(x))\,g'(x)
\]
\end{nukta}

\begin{sawal}{سوال \textenglish{20}}
حد نکالیے: \quad
\textbf{(الف)} \(\displaystyle\lim_{x\to 3}\frac{x^{4}-81}{x^{2}-x-6}\) \qquad
\textbf{(ب)} \(\displaystyle\lim_{x\to \infty}\frac{4-3x^{3}}{x^{3}-1}\)
\end{sawal}

\begin{hal}
\textbf{(الف)} تعویض سے \(\tfrac{81-81}{9-3-6}=\tfrac00\) آتا ہے۔ دونوں کی
تحلیل کیجیے --- اوپر دو مربعوں کا فرق دو بار لگتا ہے:
\[
x^{4}-81=(x^{2}-9)(x^{2}+9)=(x-3)(x+3)(x^{2}+9),
\qquad
x^{2}-x-6=(x-3)(x+2)
\]
مشترک عامل \((x-3)\) منسوخ ہو جاتا ہے:
\[
\lim_{x\to3}\frac{(x+3)(x^{2}+9)}{x+2}=\frac{(6)(18)}{5}=\frac{108}{5}=21.6
\]

\medskip
\textbf{(ب)} \(x\to\infty\) پر یہ \(\tfrac{\infty}{\infty}\) کی صورت ہے۔ اوپر
اور نیچے دونوں کو \(x^{3}\) سے تقسیم کیجیے:
\[
\lim_{x\to\infty}\frac{\tfrac{4}{x^{3}}-3}{1-\tfrac{1}{x^{3}}}
=\frac{0-3}{1-0}=-3
\]
مختصر راستہ: برابر درجے والے کثیر رقمی کسر کی حد سرکردہ عاملوں کا تناسب ہوتی
ہے، یعنی \(\tfrac{-3}{1}\)۔

\jawab{(الف) \(\tfrac{108}{5}=21.6\) \quad (ب) \(-3\)}
\end{hal}

\begin{sawal}{سوال \textenglish{21}}
\(k\) کی وہ قدر نکالیے جس کے لیے \(\displaystyle\lim_{x\to 2}f(x)\) موجود ہو،
جہاں
\[
f(x)=\begin{cases} 2-x, & x<2\\[2pt] x^{3}+k(x+1), & x\ge 2\end{cases}
\]
\end{sawal}

\begin{hal}
حد تبھی موجود ہوتی ہے جب دونوں یک طرفہ حدیں برابر ہوں:
\[
\lim_{x\to2^{-}}f(x)=2-2=0,
\qquad
\lim_{x\to2^{+}}f(x)=8+3k
\]
برابر رکھیے:
\[
8+3k=0 \;\Longrightarrow\; k=-\frac83
\]
\parakh{\(k=-\tfrac83\) رکھنے پر \(f(2^{+})=8-8=0=f(2^{-})\)۔ حد موجود ہے اور
\(0\) کے برابر ہے --- بلکہ \(f(2)=0\) بھی ہے، اس لیے فعل وہاں متصل بھی ہو
جاتا ہے۔}

\jawab{\(k=-\tfrac83\)}
\end{hal}

\begin{sawal}{سوال \textenglish{22}}
\(x\) اکائیوں کی لاگت \(C(x)=x^{2}-5x+80\) ہے۔
\textbf{(الف)} \textenglish{10} اور \textenglish{15} کے درمیان تبدیلی کی اوسط
شرح نکالیے۔
\textbf{(ب)} \textenglish{12} اکائیوں پر تبدیلی کی آنی شرح نکالیے۔
\end{sawal}

\begin{hal}
\textbf{(الف)} اوسط شرح دو نقطوں کو ملانے والے وتر کا میلان ہے:
\[
C(15)=225-75+80=230,\qquad C(10)=100-50+80=130
\]
\[
\frac{\Delta C}{\Delta x}=\frac{230-130}{15-10}=\frac{100}{5}=20
\]

\textbf{(ب)} آنی شرح مشتق ہے، یعنی ایک ہی نقطے پر مماس کا میلان۔ یہی
\emph{حاشیائی لاگت} کہلاتی ہے:
\[
C'(x)=2x-5 \;\Longrightarrow\; C'(12)=24-5=19
\]

\jawab{(الف) \(20\) فی اکائی \quad (ب) \(C'(12)=19\) فی اکائی}
\end{hal}

\begin{ghalti}
''اوسط شرح'' کے لیے دو نقطے اور تفریق چاہیے؛ ''آنی شرح'' کے لیے ایک نقطہ اور
مشتق۔ یہ الگ الگ مقداریں ہیں اور یہاں الگ جواب دیتی ہیں (\(20\) بمقابلہ
\(19\))۔ حصہ (الف) میں تفرق کرنا یا حصہ (ب) میں \(C\) میں تعویض کرنا غلط سوال کا
جواب ہے۔
\end{ghalti}

%---------------------------------------------------------------------
\section*{یونٹ \textenglish{8}: جزوی مشتق}
\markboth{یونٹ 8: جزوی مشتق}{}

\begin{nukta}[دوسرے متغیر کو ثابت رکھیے]
\(f_x\) نکالنے کے لیے \(x\) کے حوالے سے تفرق کیجیے اور \(y\) کو ایک
\emph{عدد} سمجھیے۔ \(f_y\) کے لیے \(x\) کو عدد سمجھیے۔ تفرق کے تمام معمول کے
اصول جوں کے توں لاگو رہتے ہیں۔

\medskip
\textbf{دوسرے مشتق کی جانچ} (\(f_x=f_y=0\) والے نقطے پر)، جہاں
\(D=f_{xx}f_{yy}-(f_{xy})^{2}\):
\begin{center}
\begin{tabular}{@{}p{4.4cm} p{4.4cm} p{4.0cm}@{}}
\toprule
\(D>0\) اور \(f_{xx}>0\) & \(D>0\) اور \(f_{xx}<0\) & \(D<0\) \\
اصغر (کم سے کم) & اعظم (زیادہ سے زیادہ) & زینی نقطہ \\
\bottomrule
\end{tabular}
\end{center}
\(D=0\) ہو تو جانچ کوئی فیصلہ نہیں دیتی۔
\end{nukta}

\begin{sawal}{سوال \textenglish{23}}
اگر \(z=x^{2}-y^{2}+2xy\) ہو تو ثابت کیجیے کہ
\(\pdd{z}{x}+\pdd{z}{y}=0\)۔
\end{sawal}

\begin{hal}
\(x\) کے حوالے سے (\(y\) ثابت، اس لیے \(-y^{2}\) ختم اور \(2xy\) سے \(2y\)):
\[
\pd{z}{x}=2x+2y \;\Longrightarrow\; \pdd{z}{x}=2
\]
\(y\) کے حوالے سے (\(x\) ثابت، اس لیے \(x^{2}\) ختم اور \(2xy\) سے \(2x\)):
\[
\pd{z}{y}=-2y+2x \;\Longrightarrow\; \pdd{z}{y}=-2
\]
جمع کیجیے:
\[
\pdd{z}{x}+\pdd{z}{y}=2+(-2)=0
\]
ایسا فعل \emph{ہم آہنگ} (\textenglish{harmonic}) کہلاتا ہے۔ غور کیجیے کہ
عبوری رقم \(2xy\) دونوں دوسرے مشتقوں میں کچھ نہیں ڈالتی۔

\jawab{\(z_{xx}=2\) اور \(z_{yy}=-2\)، لہٰذا \(z_{xx}+z_{yy}=0\)۔}
\end{hal}

\begin{sawal}{سوال \textenglish{24}}
ایک کمپنی کا منافع \(P(x,y)=1200+80x-2x^{2}+100y-y^{2}\) ہے، جہاں \(x\) محنت
اور \(y\) اشیا کی فی اکائی لاگت ہے۔ منافع کو زیادہ سے زیادہ کرنے والے \(x\) اور
\(y\) اور زیادہ سے زیادہ منافع نکالیے۔
\end{sawal}

\begin{hal}
یہاں \(xy\) کی رقم نہیں، اس لیے دونوں متغیر الگ الگ بہتر ہوتے ہیں:
\[
\pd{P}{x}=80-4x=0 \;\Longrightarrow\; x=20,
\qquad
\pd{P}{y}=100-2y=0 \;\Longrightarrow\; y=50
\]
\textbf{دوسرے مشتق کی جانچ:}
\[
P_{xx}=-4,\quad P_{yy}=-2,\quad P_{xy}=0,
\qquad
D=(-4)(-2)-0=8>0,\quad P_{xx}<0
\]
یعنی یہ اعظم (زیادہ سے زیادہ) ہے۔
\[
P(20,50)=1200+1600-800+5000-2500=4500
\]
\parakh{کسی بھی متغیر کو ہلانے سے منافع گھٹتا ہے:
\(P(21,50)=4498\) اور \(P(20,51)=4499\)۔}

\jawab{\(x=20\)، \(y=50\)؛ زیادہ سے زیادہ منافع \textbf{\textenglish{4{,}500}}۔}
\end{hal}

\begin{ghalti}
صرف \(P_x=0\) اور \(P_y=0\) لکھ کر ''اعظم'' کہہ دینا کافی نہیں۔ دوسرے مشتق کی
جانچ ہی اعظم، اصغر اور زینی نقطے میں فرق کرتی ہے، اور اس کے نمبر الگ ملتے ہیں۔
\end{ghalti}

%---------------------------------------------------------------------
\section*{یونٹ \textenglish{9}: بہتر بندی}
\markboth{یونٹ 9: بہتر بندی}{}

\begin{nukta}[کاروبار کے چار فعل]
\[
R(x)=x\cdot p(x),\qquad
P(x)=R(x)-C(x),\qquad
\bar{C}(x)=\frac{C(x)}{x}
\]
ان کے حاشیائی روپ مشتق ہیں: \(MR=R'\)، \(MC=C'\)۔ منافع وہاں زیادہ سے زیادہ
ہوتا ہے جہاں \(P'=0\) یعنی \(MR=MC\) ہو، بشرطیکہ \(P''<0\) ہو۔
\end{nukta}

\begin{sawal}{سوال \textenglish{25}}
\(x\) اکائیوں کی کل لاگت \(c=0.03x^{2}+8x+300\) ہے۔ کس پیداوار پر فی اکائی
اوسط لاگت کم سے کم ہو گی؟
\end{sawal}

\begin{hal}
اوسط لاگت کل لاگت تقسیم پیداوار ہے:
\[
\bar{c}(x)=\frac{0.03x^{2}+8x+300}{x}=0.03x+8+\frac{300}{x}
\]
تفرق کر کے صفر رکھیے:
\[
\bar{c}\,'(x)=0.03-\frac{300}{x^{2}}=0
\;\Longrightarrow\;
x^{2}=\frac{300}{0.03}=10{,}000
\;\Longrightarrow\;
x=100
\]
(مثبت جڑ، کیونکہ پیداوار منفی نہیں ہو سکتی۔)
\[
\bar{c}\,''(x)=\frac{600}{x^{3}}>0 \quad\text{لہٰذا یہ اصغر ہے}
\]
\[
\bar{c}(100)=3+8+3=14
\]
\parakh{\(\bar{c}(90)=14.03\) اور \(\bar{c}(110)=14.03\) --- دونوں \(14\) سے
زیادہ ہیں۔ نیز اصغر پر حاشیائی لاگت اوسط لاگت کے برابر ہوتی ہے:
\(c'(100)=0.06(100)+8=14\)۔ یہ ہمیشہ سچ ہے اور بہترین پڑتال ہے۔}

\jawab{\textbf{\textenglish{100} اکائیوں} پر اوسط لاگت کم سے کم ہوتی ہے، اور
وہاں اوسط لاگت \textbf{\textenglish{\$14}} فی اکائی ہے۔}
\end{hal}

\begin{sawal}{سوال \textenglish{26}}
طلب \(p=52-0.02x\) اور لاگت \(c(x)=400+40x\) ہے۔ کس پیداوار پر منافع زیادہ سے
زیادہ ہو گا؟ اس وقت قیمت کیا ہو گی اور منافع کتنا؟
\end{sawal}

\begin{hal}
آمدنی قیمت ضرب مقدار ہے:
\[
R(x)=x(52-0.02x)=52x-0.02x^{2}
\]
\[
P(x)=R(x)-c(x)=52x-0.02x^{2}-400-40x=12x-0.02x^{2}-400
\]
\[
P'(x)=12-0.04x=0 \;\Longrightarrow\; x=300,
\qquad
P''(x)=-0.04<0 \ \text{(اعظم)}
\]
\[
p=52-0.02(300)=46,
\qquad
P(300)=3600-1800-400=1400
\]
\parakh{\(MR=MC\) سے بھی: \(52-0.04x=40\) دیتا ہے \(x=300\) --- وہی پیداوار،
مگر معاشی اصول سے۔}

\jawab{پیداوار \textbf{\textenglish{300} اکائیاں}، قیمت
\textbf{\textenglish{\$46}} فی اکائی، اور زیادہ سے زیادہ منافع
\textbf{\textenglish{\$1{,}400}}۔}
\end{hal}

\begin{ghalti}
زیادہ سے زیادہ منافع اور زیادہ سے زیادہ آمدنی ایک چیز نہیں۔ یہاں آمدنی
\(x=1300\) پر عروج پر ہوتی ہے (جہاں \(52-0.04x=0\))، جو منافع بخش
\textenglish{300} سے بہت آگے ہے۔ اسی فعل کو بہتر کیجیے جس کا نام سوال میں ہو۔
\end{ghalti}

\begin{sawal}{سوال \textenglish{27}}
ایک کیبل کمپنی کے \textenglish{500} صارفین ہیں جو ماہانہ \textenglish{\$20}
دیتے ہیں۔ ہر \textenglish{\$0.60} کمی پر \textenglish{200} نئے صارفین آتے ہیں۔
کس نرخ پر آمدنی زیادہ سے زیادہ ہو گی اور وہ آمدنی کتنی ہو گی؟
\end{sawal}

\begin{hal}
فرض کیجیے \(n\) دفعہ \textenglish{\$0.60} کی کمی کی گئی۔ تب
\[
\text{نرخ}=20-0.6n,
\qquad
\text{صارفین}=500+200n
\]
\[
R(n)=(20-0.6n)(500+200n)=10000+3700n-120n^{2}
\]
\[
R'(n)=3700-240n=0 \;\Longrightarrow\; n=\frac{3700}{240}=15.4167,
\qquad
R''=-240<0 \ \text{(اعظم)}
\]
\[
\text{نرخ}=20-0.6(15.4167)=10.75,
\qquad
R=38{,}520.83
\]

\medskip
\textbf{مکمل عدد والا جواب۔} \(n\) پوری تعداد میں ہونا چاہیے اور صارفین بھی
پورے۔ دونوں قریبی قدریں آزمائیے:

\medskip
\begin{center}
\begin{ltrtab}
\begin{tabular}{@{}c c c c@{}}
\toprule
$n$ & نرخ & صارفین & آمدنی \\
\midrule
15 & \$11.00 & 3{,}500 & \$38{,}500 \\
16 & \$10.40 & 3{,}700 & \$38{,}480 \\
\bottomrule
\end{tabular}
\end{ltrtab}
\end{center}

\jawab{متصل (\textenglish{continuous}) حل: نرخ \textbf{\textenglish{\$10.75}}
اور آمدنی \textbf{\textenglish{\$38{,}520.83}}۔ اگر کمی پوری
\textenglish{\$0.60} کے حساب سے ہو تو بہترین نرخ \textbf{\textenglish{\$11.00}}
ماہانہ اور آمدنی \textbf{\textenglish{\$38{,}500}}۔}
\end{hal}

\begin{ghalti}
متغیر کی تعریف لکھ کر شروع کیجیے --- ''فرض کیجیے \(n\) دفعہ
\textenglish{\$0.60} کی کمی کی گئی'' --- اور نرخ اور صارفین دونوں کو اسی کے
حوالے سے لکھیے، پھر ضرب دیجیے۔ اس سوال میں زیادہ تر نمبر تفرق شروع ہونے سے
\emph{پہلے} ضائع ہوتے ہیں۔
\end{ghalti}
